% Fichier corrigé pour une compilation autonome sous MiKTeX.
% Source : SLCI/SLCI-11-DiagrammeBode/510_01_Divers/510_01_Divers.tex
\begin{corrigebox}[Corrigé extrait de la banque]

\CorrigeQuestion{1}
\textbf{Tracé asymptotique}

\begin{marginfigure}

\includegraphics[width=\linewidth]{tab_01}
\end{marginfigure}

\textbf{Positionnement du diagramme de gain}
Lorsque $\omega$ tend vers 0, le gain tend vers
$20\log 15=\SI{23,5}{dB}$.

Le pôle est situé à la pulsation de coupure
\[
\omega_c=\frac{1}{10}=\SI{0,1}{rad.s^{-1}}.
\]
Le diagramme de gain possède donc une asymptote horizontale à
\SI{23,5}{dB} avant $\omega_c$, puis une pente de
\SI{-20}{dB/decade}. Le déphasage évolue de $0^\circ$ vers
$-90^\circ$ et vaut $-45^\circ$ à $\omega_c$.

\begin{figure}[!h]
\centering
\includegraphics[width=\linewidth]{bode_01}
\end{figure}

\CorrigeQuestion{2}
Pour une période de \SI{60}{s}, la pulsation est
\[
\omega=\frac{2\pi}{T}\simeq\SI{0,1}{rad.s^{-1}}.
\]
À cette pulsation, le gain est de \SI{20}{dB} et le déphasage de
$-\dfrac{\pi}{4}$.

On a donc $20\log(S/E)=20$, soit $S=10E$. Le signal d'entrée est
\[
e(t)=10\sin(0,1t)
\]
et le signal de sortie est
\[
s(t)=100\sin\left(0,1t-\frac{\pi}{4}\right).
\]

\end{corrigebox}
